% !TEX root = ../../thesis.tex

\section{Towards Multiplicity One}\label{sec:multiplicity-one}

The sheet counting argument in Proposition~\ref{prop:integrality-nested} yields a parity constraint on the density $\theta$ of the limit varifold $V$. At pure cancellation points ($p \in \spt\|V'\| \setminus \spt\|\partial^*\Omega(x_0)\|$), the cancelling sheets pair up, giving $\theta(p) = 2m$ for some $m \geq 1$. At mixed regular points ($p \in \mathrm{reg}(\partial^*\Omega(x_0)) \cap \spt\|V'\|$), one tracking sheet persists and the rest pair, giving $\theta(p) = 1 + 2m$. Thus multiplicity one ($\theta \equiv 1$) is equivalent to $m = 0$ everywhere, which is precisely the no-mass-cancellation condition (Definition~\ref{def:mass-cancellation}).

A natural strategy for proving multiplicity one is the following. Suppose $V = k|\Sigma|$ with $k \geq 2$ and $\Sigma$ a smooth closed embedded minimal hypersurface (as guaranteed by Theorem~\ref{thm:non-excessive-main}). One attempts to construct an $I$-replacement family (Definition~\ref{def:excessive-interval}) on an interval $I \ni x_0$ by performing surgery on near-maximal slices: for $x$ with $\mathrm{Per}(\Omega(x))$ close to $W = k \cdot \mathcal{H}^n(\Sigma)$, replace $\Omega(x)$ inside a tubular neighborhood $B_\varepsilon(\Sigma)$ by $\Omega(x_0)$, removing the $k - 1$ cancelling sheet pairs. Since $\mathrm{Per}(\Omega(x_0)) \leq \mathcal{H}^n(\Sigma) = W/k$, and the cutting cost on $\partial B_\varepsilon(\Sigma)$ vanishes as $x \to x_0$ by $L^1$ convergence, the modified perimeter satisfies $\mathrm{Per}(\tilde{\Omega}(x)) \approx W/k \ll W$. This would make $x_0$ excessive, contradicting non-excessiveness.

The obstruction is that $L^1$ convergence of Caccioppoli sets does not control the spatial distribution of perimeter. For the specific min-max sequence $x_i \to x_0$, varifold convergence $|\partial^*\Omega(x_i)| \to k|\Sigma|$ guarantees that the perimeter of $\Omega(x_i)$ concentrates in $B_\varepsilon(\Sigma)$, so the surgery reduces perimeter as claimed. But for a general parameter $x$ near $x_0$, the boundary $\partial^*\Omega(x)$ need not be near $\Sigma$: a slice with $\mathrm{Per}(\Omega(x)) \approx (k-1)\mathcal{H}^n(\Sigma)$ and $\mathrm{Per}(\Omega(x), B_\varepsilon(\Sigma)) \approx 0$ would see its perimeter \emph{increase} to $\approx (k-1)\mathcal{H}^n(\Sigma) + \mathcal{H}^n(\Sigma) = W$ after the replacement, violating the mass bound required of an $I$-replacement family. The $I$-replacement definition demands $\limsup$ mass $< W$ for \emph{all} $x \in I$, and one cannot restrict the surgery to the subsequence where it is effective.

This gap reflects a general difficulty: the non-excessive framework controls the limit varifold through one-sided homotopic minimizing properties of critical slices, but does not directly constrain the geometry of nearby non-critical slices. Zhou's proof of multiplicity one~\cite{Zho20} overcomes analogous obstructions using the substantially deeper machinery of prescribed mean curvature (PMC) min-max theory, which provides finer control on the interaction between the sweepout parameter and the geometry of individual slices.

\begin{question}\label{q:mult-one}
Can multiplicity one be established within the CLS non-excessive framework, using nestedness and the sheet decomposition (Lemma~\ref{lem:sheet-decomposition}) to rule out mass cancellation?
\end{question}
